\documentclass[11pt]{article}
% ===========================================================================
%  Shared preamble for all MPM2D worksheets — input right after \documentclass.
%  (Files beginning with "_" are skipped by mpm2d-worksheet-publish.mjs.)
% ===========================================================================
\usepackage[letterpaper,top=2.2cm,bottom=1.9cm,left=1.6cm,right=1.6cm,headheight=28pt,footskip=24pt]{geometry}
\usepackage{amsmath,amssymb}
\usepackage[T1]{fontenc}
\usepackage[utf8]{inputenc}
\usepackage{lmodern}
\usepackage{xcolor}
\usepackage[most]{tcolorbox}
\usepackage{enumitem}
\usepackage{multicol}
\usepackage{fancyhdr}
\usepackage{tikz}
\usetikzlibrary{calc}
\usepackage{pgfplots}
\pgfplotsset{compat=1.18}

\definecolor{exblue}{HTML}{4A90E2}
\definecolor{exbg}{HTML}{E6F3FF}
\definecolor{qorange}{HTML}{E69138}
\definecolor{qbg}{HTML}{FFF7CC}
\definecolor{keygray}{HTML}{475569}
\definecolor{keybg}{HTML}{F1F5F9}
\definecolor{iagreen}{HTML}{14653B}
\definecolor{titleblue}{HTML}{1D4ED8}
% Rotating palette so each worked example gets its own colour.
\definecolor{ex0}{HTML}{2563A0}\definecolor{ex1}{HTML}{3B7D3B}\definecolor{ex2}{HTML}{8A5A00}
\definecolor{ex3}{HTML}{A3327A}\definecolor{ex4}{HTML}{6D28A3}\definecolor{ex5}{HTML}{B08900}
\definecolor{ex6}{HTML}{0E7490}\definecolor{ex7}{HTML}{9A3412}\definecolor{ex8}{HTML}{15803D}

\pagestyle{fancy}
\fancyhf{}
\fancyhead[L]{\footnotesize NAME: \underline{\hspace{3.4cm}}}
\fancyhead[C]{\textbf{Integration Academy}\\[-2pt]{\scriptsize Math Simplified \textperiodcentered{} Success Amplified}}
\fancyhead[R]{\footnotesize MPM2D}
\fancyfoot[L]{\scriptsize dr.merie@IntegrationAcademy.ca}
\fancyfoot[C]{\scriptsize\thepage}
\fancyfoot[R]{\scriptsize IntegrationAcademy.ca}
\renewcommand{\headrulewidth}{1.2pt}
\renewcommand{\footrulewidth}{0.4pt}
\setlength{\parindent}{0pt}
\setlength{\parskip}{4pt}

% Examples are NOT breakable, so each example + its solution stays on one page.
\newtcolorbox{example}[2]{colframe=#1,colback=#1!7,coltitle=white,colbacktitle=#1,fonttitle=\bfseries,title={#2},arc=3pt,boxrule=1pt}
\newtcolorbox{qbox}[1]{colframe=qorange,colback=qbg,coltitle=white,colbacktitle=qorange,fonttitle=\bfseries,title={#1},arc=3pt,breakable,boxrule=1pt}
\newtcolorbox{keybox}{colframe=keygray,colback=keybg,coltitle=white,colbacktitle=keygray,fonttitle=\bfseries,title={$\checkmark$\ Answer Key},arc=3pt,breakable,boxrule=1pt}
\newtcolorbox{recapbox}{colframe=keygray,colback=keybg,coltitle=white,colbacktitle=keygray,fonttitle=\bfseries,title={Question Recap},arc=3pt,breakable,boxrule=1pt}
\newtcolorbox{ideabox}{colframe=titleblue,colback=titleblue!6,coltitle=white,colbacktitle=titleblue,fonttitle=\bfseries,title={The Big Ideas},arc=3pt,breakable,boxrule=1.2pt}
\newtcolorbox{learnbox}{colframe=iagreen,colback=iagreen!5,coltitle=white,colbacktitle=iagreen,fonttitle=\bfseries,title={Core Concepts},arc=3pt,breakable,boxrule=1.2pt}
\newcommand{\lh}[1]{\par\smallskip{\bfseries\color{iagreen}#1}\par\nobreak\smallskip}

\newcommand{\soln}{\par\smallskip\textit{\textbf{Solution.}}\ }
\newcommand{\workspace}[1]{\par\vspace{3pt}\fcolorbox{black!30}{white}{\begin{minipage}[t][#1][t]{\dimexpr\linewidth-2\fboxsep\relax}\ \end{minipage}}\par}
\newcommand{\sech}[1]{\par\vspace{8pt}{\Large\bfseries\color{iagreen}#1}\par\nobreak\vspace{2pt}{\color{black!20}\hrule height1pt}\vspace{6pt}}

% A compact coordinate plot sized for an example box: \eplot{xmin}{xmax}{ymin}{ymax}{body}
\newcommand{\eplot}[5]{\begin{center}\begin{tikzpicture}
\begin{axis}[width=6.8cm,height=4.4cm,axis lines=middle,xlabel={\tiny$x$},ylabel={\tiny$y$},
grid=both,grid style={gray!16},xmin=#1,xmax=#2,ymin=#3,ymax=#4,
xtick distance=2,ytick distance=2,tick label style={font=\tiny},axis line style={-},samples=120]
#5
\end{axis}\end{tikzpicture}\end{center}}

% Same as \eplot but with its own tick spacing: \eplotx{xmin}{xmax}{ymin}{ymax}{xtick}{ytick}{body}
\newcommand{\eplotx}[7]{\begin{center}\begin{tikzpicture}
\begin{axis}[width=8.4cm,height=5.4cm,axis lines=middle,xlabel={\tiny$x$},ylabel={\tiny$y$},
grid=both,grid style={gray!16},xmin=#1,xmax=#2,ymin=#3,ymax=#4,
xtick distance=#5,ytick distance=#6,tick label style={font=\tiny},axis line style={-},samples=120]
#7
\end{axis}\end{tikzpicture}\end{center}}

% A sine-curve plot (degrees on x, 0..360): \splot{ymin}{ymax}{body}
\newcommand{\splot}[3]{\begin{center}\begin{tikzpicture}
\begin{axis}[width=8.6cm,height=4.6cm,axis lines=middle,xlabel={\tiny$x\,(^\circ)$},ylabel={\tiny$y$},
grid=both,grid style={gray!16},xmin=0,xmax=360,ymin=#1,ymax=#2,
xtick distance=90,ytick distance=1,tick label style={font=\tiny},axis line style={-},samples=180]
#3
\end{axis}\end{tikzpicture}\end{center}}

% A small coordinate plot: \plot{xmin}{xmax}{ymin}{ymax}{body}
\newcommand{\plot}[5]{\begin{center}\begin{tikzpicture}
\begin{axis}[width=8.4cm,height=6cm,axis lines=middle,xlabel={\scriptsize$x$},ylabel={\scriptsize$y$},
grid=both,grid style={gray!16},xmin=#1,xmax=#2,ymin=#3,ymax=#4,
xtick distance=2,ytick distance=2,tick label style={font=\tiny},axis line style={-}]
#5
\end{axis}\end{tikzpicture}\end{center}}

% Right triangle (angle theta at bottom-left, right angle at bottom-right):
%   \rtri{drawBase}{drawHeight}{oppLabel}{adjLabel}{hypLabel}
\newcommand{\rtri}[5]{\begin{center}\begin{tikzpicture}[scale=0.85]
\coordinate (A) at (0,0);\coordinate (B) at (#1,0);\coordinate (C) at (#1,#2);
\draw[thick,exblue] (A)--(B)--(C)--cycle;
\draw[black] ($(B)+(-0.32,0)$)--($(B)+(-0.32,0.32)$)--($(B)+(0,0.32)$);
\node[below] at ($(A)!0.5!(B)$){#4};
\node[right] at ($(B)!0.5!(C)$){#3};
\node[above left] at ($(A)!0.5!(C)$){#5};
\node at (0.7,0.22){$\theta$};
\end{tikzpicture}\end{center}}

% General (oblique) triangle with vertex labels and opposite-side labels:
%   \gtri{Alabel}{Blabel}{Clabel}{aSide}{bSide}{cSide}
\newcommand{\gtri}[6]{\begin{center}\begin{tikzpicture}[scale=0.85]
\coordinate (A) at (0,0);\coordinate (B) at (5,0);\coordinate (C) at (1.6,3);
\draw[thick,exblue] (A)--(B)--(C)--cycle;
\node[below left] at (A){#1};\node[below right] at (B){#2};\node[above] at (C){#3};
\node[below] at ($(A)!0.5!(B)$){#6};
\node[right] at ($(B)!0.5!(C)$){#4};
\node[left] at ($(A)!0.5!(C)$){#5};
\end{tikzpicture}\end{center}}

\fancyhead[R]{\footnotesize MHF4U \textperiodcentered{} 2.1}
\begin{document}

\begin{center}
{\LARGE\bfseries\color{titleblue}Worksheet 2.1 --- Dividing Polynomials}\\[2pt]
{\small Grade 12 Advanced Functions (MHF4U) \textperiodcentered{} Unit 2: Polynomial Equations and Inequalities}
\end{center}

\begin{ideabox}
Long division always works; synthetic division is a shortcut for dividing by $x-a$. Every result is a division statement $P(x)=D(x)Q(x)+R(x)$.
\begin{itemize}[leftmargin=*,itemsep=2pt]
  \item Long division: divide, multiply, subtract, bring down.
  \item Synthetic (by $x-a$): bring down, multiply by $a$, add.
  \item $\deg R<\deg D$; remainder $0$ means $D$ is a factor.
\end{itemize}
\end{ideabox}

\begin{learnbox}
\lh{The division identity}Polynomial division gives $P(x)=D(x)Q(x)+R$, with the remainder $R$ lower in degree than the divisor $D$.
\lh{Long and synthetic}Long division works for any divisor; \textbf{synthetic division} is a fast shortcut when the divisor is $x-a$.
\lh{Fill the gaps}Insert zero coefficients for missing powers so the terms line up correctly through the division.
\end{learnbox}

\sech{Worked Examples}

\begin{example}{ex0}{Example 1 --- Exact long division}
$(x^2+8x+15)\div(x+3)$.\soln $x^2\div x=x$; subtract $x^2+3x$ → $5x+15$; $5x\div x=5$; subtract → $0$. Quotient $x+5$, R $0$. The graph confirms a zero at $x=-3$:\eplot{-7}{0}{-2}{6}{\addplot[exblue,very thick,domain=-6.4:-1.6]{x^2+8*x+15};\addplot[red,only marks,mark=*,mark size=1.6pt] coordinates {(-3,0) (-5,0)};}\begin{center}\renewcommand{\arraystretch}{1.2}\begin{tabular}{cccc}
&  & $x$ & $+5$ \\
\cline{2-4}
\multicolumn{1}{r|}{$x+3$} & $x^2$ & $+8x$ & $+15$ \\
\multicolumn{1}{r|}{} & $x^2$ & $+3x$ &  \\
\cline{2-4}
\multicolumn{1}{r|}{} &  & $+5x$ & $+15$ \\
\multicolumn{1}{r|}{} &  & $+5x$ & $+15$ \\
\cline{2-4}
\multicolumn{1}{r|}{} &  &  & $0$ \\
\end{tabular}\end{center}
\end{example}

\begin{example}{ex1}{Example 2 --- Exact long division}
$(x^2+7x+10)\div(x+5)$.\soln Quotient $x+2$, R $0$, so $x^2+7x+10=(x+5)(x+2)$.\begin{center}\renewcommand{\arraystretch}{1.2}\begin{tabular}{cccc}
&  & $x$ & $+2$ \\
\cline{2-4}
\multicolumn{1}{r|}{$x+5$} & $x^2$ & $+7x$ & $+10$ \\
\multicolumn{1}{r|}{} & $x^2$ & $+5x$ &  \\
\cline{2-4}
\multicolumn{1}{r|}{} &  & $+2x$ & $+10$ \\
\multicolumn{1}{r|}{} &  & $+2x$ & $+10$ \\
\cline{2-4}
\multicolumn{1}{r|}{} &  &  & $0$ \\
\end{tabular}\end{center}
\end{example}

\begin{example}{ex2}{Example 3 --- Synthetic division}
$(x^3-5x^2+2x+8)\div(x-4)$.\soln Coefficients $1,-5,2,8$, root $4$: bring down $1$; $1{\cdot}4=4,\ -5{+}4=-1$; $-1{\cdot}4=-4,\ 2{-}4=-2$; $-2{\cdot}4=-8,\ 8{-}8=0$. Quotient $x^2-x-2$, R $0$.\begin{center}\renewcommand{\arraystretch}{1.2}\begin{tabular}{ccccc}
\multicolumn{1}{r|}{$4$} & $1$ & $-5$ & $2$ & $8$ \\
\multicolumn{1}{r|}{} &  & $4$ & $-4$ & $-8$ \\
\cline{2-5}
\multicolumn{1}{r|}{} & $1$ & $-1$ & $-2$ & $0$ \\
\end{tabular}\end{center}
\end{example}

\begin{example}{ex3}{Example 4 --- Division statement}
Write the statement for $(x^2+2x-15)\div(x+5)$.\soln $x^2+2x-15=(x+5)(x-3)+0$.\begin{center}\renewcommand{\arraystretch}{1.2}\begin{tabular}{cccc}
&  & $x$ & $-3$ \\
\cline{2-4}
\multicolumn{1}{r|}{$x+5$} & $x^2$ & $+2x$ & $-15$ \\
\multicolumn{1}{r|}{} & $x^2$ & $+5x$ &  \\
\cline{2-4}
\multicolumn{1}{r|}{} &  & $-3x$ & $-15$ \\
\multicolumn{1}{r|}{} &  & $-3x$ & $-15$ \\
\cline{2-4}
\multicolumn{1}{r|}{} &  &  & $0$ \\
\end{tabular}\end{center}
\end{example}

\begin{example}{ex4}{Example 5 --- Nonzero remainder}
$(x^2+5x+9)\div(x+2)$.\soln Quotient $x+3$; $(x+2)(x+3)=x^2+5x+6$; subtract → $3$. So $x^2+5x+9=(x+2)(x+3)+3$.\begin{center}\renewcommand{\arraystretch}{1.2}\begin{tabular}{cccc}
&  & $x$ & $+3$ \\
\cline{2-4}
\multicolumn{1}{r|}{$x+2$} & $x^2$ & $+5x$ & $+9$ \\
\multicolumn{1}{r|}{} & $x^2$ & $+2x$ &  \\
\cline{2-4}
\multicolumn{1}{r|}{} &  & $+3x$ & $+9$ \\
\multicolumn{1}{r|}{} &  & $+3x$ & $+6$ \\
\cline{2-4}
\multicolumn{1}{r|}{} &  &  & $+3$ \\
\end{tabular}\end{center}
\end{example}

\begin{example}{ex5}{Example 6 --- Synthetic}
$(x^3+8)\div(x+2)$.\soln Coefficients $1,0,0,8$, root $-2$: $1,-2,4,0$. Quotient $x^2-2x+4$, R $0$.\begin{center}\renewcommand{\arraystretch}{1.2}\begin{tabular}{ccccc}
\multicolumn{1}{r|}{$-2$} & $1$ & $0$ & $0$ & $8$ \\
\multicolumn{1}{r|}{} &  & $-2$ & $4$ & $-8$ \\
\cline{2-5}
\multicolumn{1}{r|}{} & $1$ & $-2$ & $4$ & $0$ \\
\end{tabular}\end{center}
\end{example}

\begin{example}{ex6}{Example 7 --- Linear over linear}
$(x^2-x-12)\div(x-4)$.\soln Quotient $x+3$, R $0$, so $(x-4)(x+3)$.\begin{center}\renewcommand{\arraystretch}{1.2}\begin{tabular}{cccc}
&  & $x$ & $+3$ \\
\cline{2-4}
\multicolumn{1}{r|}{$x-4$} & $x^2$ & $-x$ & $-12$ \\
\multicolumn{1}{r|}{} & $x^2$ & $-4x$ &  \\
\cline{2-4}
\multicolumn{1}{r|}{} &  & $+3x$ & $-12$ \\
\multicolumn{1}{r|}{} &  & $+3x$ & $-12$ \\
\cline{2-4}
\multicolumn{1}{r|}{} &  &  & $0$ \\
\end{tabular}\end{center}
\end{example}

\begin{example}{ex7}{Example 8 --- Remainder}
$(x^2+6x+3)\div(x+1)$.\soln Quotient $x+5$, remainder $-2$: $x^2+6x+3=(x+1)(x+5)-2$.\begin{center}\renewcommand{\arraystretch}{1.2}\begin{tabular}{cccc}
&  & $x$ & $+5$ \\
\cline{2-4}
\multicolumn{1}{r|}{$x+1$} & $x^2$ & $+6x$ & $+3$ \\
\multicolumn{1}{r|}{} & $x^2$ & $+x$ &  \\
\cline{2-4}
\multicolumn{1}{r|}{} &  & $+5x$ & $+3$ \\
\multicolumn{1}{r|}{} &  & $+5x$ & $+5$ \\
\cline{2-4}
\multicolumn{1}{r|}{} &  &  & $-2$ \\
\end{tabular}\end{center}
\end{example}

\begin{example}{ex8}{Example 9 --- Cubic remainder}
$(x^3+3x^2-2)\div(x-1)$.\soln Synthetic with root $1$ on $1,3,0,-2$: $1,4,4,2$. Quotient $x^2+4x+4$, R $2$.\begin{center}\renewcommand{\arraystretch}{1.2}\begin{tabular}{ccccc}
\multicolumn{1}{r|}{$1$} & $1$ & $3$ & $0$ & $-2$ \\
\multicolumn{1}{r|}{} &  & $1$ & $4$ & $4$ \\
\cline{2-5}
\multicolumn{1}{r|}{} & $1$ & $4$ & $4$ & $2$ \\
\end{tabular}\end{center}
\end{example}

\clearpage
\sech{Your Turn --- Show Your Work}
For each question, show all steps clearly in the space provided.

\begin{qbox}{Question 1}
$(x^2+10x+21)\div(x+3)$?\workspace{2.4cm}
\end{qbox}

\begin{qbox}{Question 2}
$(x^2-3x-10)\div(x-5)$?\workspace{2.4cm}
\end{qbox}

\begin{qbox}{Question 3}
Synthetic: $(x^3-27)\div(x-3)$?\workspace{2.4cm}
\end{qbox}

\begin{qbox}{Question 4}
$(x^2+3x+7)\div(x+2)$, with remainder?\workspace{2.4cm}
\end{qbox}

\begin{qbox}{Question 5}
Write the division statement for $(x^2+8x+5)\div(x+3)$.\workspace{2.4cm}
\end{qbox}

\begin{qbox}{Question 6}
$(x^2+9x+20)\div(x+4)$?\workspace{2.4cm}
\end{qbox}

\begin{qbox}{Question 7}
$(x^3-8)\div(x-2)$?\workspace{2.4cm}
\end{qbox}

\begin{qbox}{Question 8}
$(x^2-5x+6)\div(x-3)$?\workspace{2.4cm}
\end{qbox}

\begin{qbox}{Question 9}
$(x^3+4x^2-7)\div(x-2)$, quotient and remainder?\workspace{2.4cm}
\end{qbox}

\begin{qbox}{Question 10}
Is $x+5$ a factor of $x^2+3x-10$? (use the remainder)\workspace{2.4cm}
\end{qbox}

\begin{qbox}{Question 11}
$(2x^2+7x+3)\div(x+3)$?\workspace{2.4cm}
\end{qbox}

\begin{qbox}{Question 12}
$(x^2+1)\div(x+1)$, remainder?\workspace{2.4cm}
\end{qbox}

\begin{qbox}{Question 13 --- Challenge}
Divide $(x^3-2x^2-5x+6)\div(x-1)$ and state quotient and remainder.\workspace{3.5cm}
\end{qbox}

\clearpage
\sech{Question Recap \& Answer Key}

\begin{recapbox}
\begin{multicols}{2}
\small
\begin{enumerate}[leftmargin=*,itemsep=3pt]
  \item $(x^2+10x+21)\div(x+3)$?
  \item $(x^2-3x-10)\div(x-5)$?
  \item Synthetic: $(x^3-27)\div(x-3)$?
  \item $(x^2+3x+7)\div(x+2)$, with remainder?
  \item Write the division statement for $(x^2+8x+5)\div(x+3)$.
  \item $(x^2+9x+20)\div(x+4)$?
  \item $(x^3-8)\div(x-2)$?
  \item $(x^2-5x+6)\div(x-3)$?
  \item $(x^3+4x^2-7)\div(x-2)$, quotient and remainder?
  \item Is $x+5$ a factor of $x^2+3x-10$? (use the remainder)
  \item $(2x^2+7x+3)\div(x+3)$?
  \item $(x^2+1)\div(x+1)$, remainder?
  \item Divide $(x^3-2x^2-5x+6)\div(x-1)$ and state quotient and remainder.
\end{enumerate}
\end{multicols}
\end{recapbox}

\begin{keybox}
\begin{multicols}{2}
\small
\begin{enumerate}[leftmargin=*,itemsep=3pt]
  \item $x+7$, R 0
  \item $x+2$, R 0
  \item $x^2+3x+9$, R 0
  \item $x+1$, R 5
  \item $(x+3)(x+5)-10$
  \item $x+5$, R 0
  \item $x^2+2x+4$, R 0
  \item $x-2$, R 0
  \item $x^2+6x+12$, R 17
  \item yes (R 0)
  \item $2x+1$, R 0
  \item $x-1$, R 2
  \item $x^2-x-6$, R 0
\end{enumerate}
\end{multicols}
\end{keybox}

\end{document}
